\documentclass[../DoAn.tex]{subfiles}
\begin{document}

Chương này trình bày chi tiết về mặt kỹ thuật của hệ thống JQK Study, bao gồm thiết kế kiến trúc phần mềm, thiết kế chi tiết cơ sở dữ liệu và các lớp nghiệp vụ, cùng kết quả xây dựng ứng dụng, phương pháp kiểm thử và mô hình triển khai thực tế.

\section{Thiết kế kiến trúc}
\label{section:4.1}

\subsection{Lựa chọn kiến trúc phần mềm}
Hệ thống JQK Study được phát triển dựa trên **kiến trúc Microservices** (kiến trúc vi dịch vụ), chia nhỏ hệ thống thành các khối dịch vụ độc lập chịu trách nhiệm cho từng miền nghiệp vụ riêng biệt. Việc này cho phép mỗi thành phần hoạt động, lưu trữ dữ liệu độc lập và giao tiếp với nhau qua mạng thông qua các giao thức gọi hàm từ xa tốc độ cao gRPC hoặc thông qua hệ thống truyền tin nhắn bất đồng bộ Kafka.

Cụ thể, hệ thống được chia làm hai tầng chính:
\begin{itemize}
    \item \textbf{Tầng Client (Frontend):} Được xây dựng bằng Next.js (chạy ở cổng `4001`), giao tiếp với Backend thông qua các truy vấn REST API gửi tới API Gateway.
    \item \textbf{Tầng Services (Backend):} Bao gồm API Gateway đóng vai trò điểm điều phối trung tâm và các vi dịch vụ nghiệp vụ nội bộ (chạy ở các cổng gRPC riêng biệt).
\end{itemize}

\subsection{Thiết kế tổng quan}
Mô hình kiến trúc tổng quan của JQK Study bao gồm các dịch vụ sau:
\begin{itemize}
    \item \textbf{API Gateway (Gin Web Framework - Port 8081):} Nhận các yêu cầu HTTP từ frontend, thực hiện phân quyền (authorization) qua JWT, kiểm tra CORS, định tuyến (routing) yêu cầu và chuyển đổi chúng thành các cuộc gọi gRPC tới các dịch vụ phía sau.
    \item \textbf{User Service (Port 50051):} Quản lý thông tin người dùng, mật khẩu (mã hóa bcrypt), phân quyền vai trò (Admin, Instructor, Student), quản lý lớp học (Class) và thành viên lớp học.
    \item \textbf{Exam Service (Port 50052):} Dịch vụ lõi quan trọng nhất, quản lý ngân hàng câu hỏi, chủ đề, phần thi, sinh đề thi tự động, xử lý quá trình học viên làm bài thi, lưu đáp án nháp, chấm điểm tự động và ghi nhận nhật ký vi phạm (blur tab, copy/paste).
    \item \textbf{Course Service (Port 50053):} Quản lý các khóa học, chương học, bài giảng video/tài liệu, đăng ký khóa học của học viên và theo dõi tiến độ hoàn thành bài học.
    \item \textbf{Notification Service (Port 50054):} Hoạt động ở chế độ bất đồng bộ, đăng ký lắng nghe các sự kiện từ hàng đợi Apache Kafka để tự động gửi thông báo qua email hoặc hệ thống cho người dùng khi có sự kiện đăng ký thành công, nộp bài thi hoặc gán lớp học mới.
\end{itemize}

Sự phụ thuộc giữa các dịch vụ tuân thủ nguyên tắc lỏng lẻo (loose coupling): các dịch vụ nghiệp vụ không trực tiếp gọi cơ sở dữ liệu của nhau mà chỉ tương tác qua gRPC client được định nghĩa trước bằng các tệp giao ước gRPC Protocol Buffers.

\begin{figure}[H]
    \centering
    \includegraphics[width=0.85\textwidth]{Hinhve/Picture1.png}
    \caption{Biểu đồ phân tầng và phụ thuộc gói của hệ thống}
    \label{fig:Fig1}
\end{figure}

\subsection{Thiết kế chi tiết gói (Package Design)}
Mỗi vi dịch vụ backend trong Go được tổ chức cấu trúc thư mục chuẩn nhằm phân tách rõ ràng các tầng trách nhiệm:
\begin{itemize}
    \item \textbf{Tầng Cmd (`cmd/`):} Điểm khởi đầu của ứng dụng, chịu trách nhiệm đọc tệp cấu hình môi trường (.env), thiết lập kết nối cơ sở dữ liệu, khởi tạo máy chủ gRPC và lắng nghe các cổng.
    \item \textbf{Tầng Domain (`internal/domain/`):} Chứa định nghĩa các mô hình dữ liệu (structs) và các interface cho tầng Service và Repository. Tầng này độc lập hoàn toàn, không phụ thuộc vào bất kỳ thư viện ngoài nào (Domain-Driven Design).
    \item \textbf{Tầng Service (`internal/service/`):} Nơi chứa logic nghiệp vụ chính (Business Logic), thực hiện kiểm tra điều kiện, xử lý luồng công việc và tương tác với tầng Repository.
    \item \textbf{Tầng Database/Repository (`internal/databases/`):} Thực hiện các câu lệnh truy vấn dữ liệu chi tiết vào cơ sở dữ liệu PostgreSQL hoặc cache Redis sử dụng thư viện GORM.
    \item \textbf{Tầng Infrastructure (`internal/infrastructure/`):} Triển khai gRPC server handler để map dữ liệu gRPC request thành domain model, đồng thời chứa Kafka consumer/producer.
\end{itemize}

\begin{figure}[H]
    \centering
    \includegraphics[width=0.85\textwidth]{Hinhve/Picture2.png}
    \caption{Biểu đồ thiết kế chi tiết lớp và mối quan hệ giữa các thành phần}
    \label{fig:Fig2}
\end{figure}

\section{Thiết kế chi tiết}
\label{section:4.2}

\subsection{Thiết kế giao diện}
Giao diện JQK Study được phát triển theo triết lý hiện đại, tương thích tốt với nhiều kích thước màn hình (Responsive Design) từ thiết bị di động đến máy tính để bàn nhờ TailwindCSS. Các quy chuẩn thiết kế giao diện bao gồm:
\begin{itemize}
    \item **Hệ màu:** Sử dụng hệ màu tối (dark mode) làm chủ đạo với các tông màu xám lạnh kết hợp màu xanh ngọc (cyan) tạo nét cao cấp.
    \item **Màn hình làm bài thi:** Được thiết kế tối giản, tập trung tối đa vào nội dung câu hỏi. Có khu vực cố định hiển thị thời gian đếm ngược trực quan và danh sách trạng thái các câu hỏi đã làm/chưa làm.
    \item **Giao diện Giảng viên:** Dashboard trực quan sử dụng thư viện biểu đồ Recharts hiển thị thống kê điểm số, phân phối phổ điểm và danh sách cảnh báo vi phạm chi tiết của từng sinh viên.
\end{itemize}

\subsection{Thiết kế lớp nghiệp vụ}
Dựa trên kiến trúc mã nguồn nghiệp vụ của các dịch vụ, các interface chính được thiết kế để điều phối hoạt động bao gồm:

\textbf{1. Phân hệ User Service:}
\begin{itemize}
    \item `UserService`: Định nghĩa các phương thức nghiệp vụ đăng ký (`Register`), đăng nhập (`Login`), lấy thông tin tài khoản (`GetProfile`), tạo lớp học (`CreateClass`), và quản lý thành viên lớp học.
    \item `UserRepository`: Cung cấp giao diện tương tác dữ liệu như lưu người dùng mới (`CreateUser`), tìm kiếm tài khoản theo email (`GetUserByEmail`), tạo lớp học (`CreateClass`), và gán học viên vào lớp.
\end{itemize}

\textbf{2. Phân hệ Course Service:}
\begin{itemize}
    \item `CourseService`: Thực hiện nghiệp vụ tạo khóa học, quản lý chương học, bài giảng video, ghi nhận tiến độ xem bài giảng và đăng ký khóa học.
    \item `CourseRepository`: Thực hiện lưu trữ thông tin khóa học (`CreateCourse`), bài học (`CreateLesson`), lưu tiến trình học tập (`CreateEnrollment`, `UpdateProgress`).
\end{itemize}

\textbf{3. Phân hệ Exam Service:}
\begin{itemize}
    \item `ExamService`: Xử lý nghiệp vụ tạo đề thi thủ công hoặc tự động (`CreateExam`, `GenerateExam`), lưu trữ câu trả lời của thí sinh theo thời gian thực (`SubmitAnswer`), chấm điểm tự động (`SubmitExam`), ghi nhận vi phạm gian lận và kết xuất kết quả thi.
    \item `ExamRepository`: Quản lý các thao tác ghi dữ liệu đề thi (`CreateExam`), ngân hàng câu hỏi (`CreateQuestion`), tạo bài làm mới (`CreateSubmission`), lưu kết quả làm bài của thí sinh (`SaveUserAnswer`), và ghi lại vi phạm (`LogViolation`).
\end{itemize}

\subsection{Thiết kế cơ sở dữ liệu}
Hệ thống sử dụng cơ sở dữ liệu quan hệ PostgreSQL với các bảng được phân tách độc lập theo từng vi dịch vụ để đảm bảo tính cô lập dữ liệu.

\textbf{1. Các bảng dữ liệu của User Service:}
\begin{itemize}
    \item \texttt{roles}: Lưu trữ vai trò người dùng (id, name).
    \item \texttt{users}: Thông tin người dùng (id, \texttt{full\_name}, email, password, \texttt{role\_id}, \texttt{created\_at}, \texttt{updated\_at}).
    \item \texttt{classes}: Lưu thông tin lớp học (id, name, code, description, \texttt{teacher\_id}, \texttt{is\_active}).
    \item \texttt{class\_members}: Liên kết nhiều-nhiều giữa lớp học và học viên (\texttt{class\_id}, \texttt{user\_id}, role, status, \texttt{joined\_at}).
\end{itemize}

\textbf{2. Các bảng dữ liệu của Course Service:}
\begin{itemize}
    \item \texttt{courses}: Lưu thông tin khóa học (id, title, description, \texttt{instructor\_id}, price, \texttt{is\_published}).
    \item \texttt{course\_sections}: Phân chia chương học (id, \texttt{course\_id}, title, \texttt{order\_index}).
    \item \texttt{lessons}: Chi tiết bài học (id, \texttt{section\_id}, title, \texttt{type\_id}, \texttt{content\_url}, \texttt{duration\_seconds}, \texttt{order\_index}).
    \item \texttt{enrollments}: Học viên tham gia khóa học (\texttt{user\_id}, \texttt{course\_id}, \texttt{enrolled\_at}).
    \item \texttt{lesson\_progresses}: Theo dõi tiến độ xem bài học (\texttt{user\_id}, \texttt{lesson\_id}, \texttt{completed\_at}).
\end{itemize}

\textbf{3. Các bảng dữ liệu của Exam Service (Dịch vụ thi):}
\begin{itemize}
    \item \texttt{topics} và \texttt{exam\_sections}: Quản lý danh mục ngân hàng câu hỏi.
    \item \texttt{questions}: Lưu trữ nội dung câu hỏi, loại câu hỏi (một đáp án/nhiều đáp án), độ khó (dễ, trung bình, khó), giải thích đáp án.
    \item \texttt{choices}: Danh sách các câu trả lời lựa chọn của câu hỏi (id, \texttt{question\_id}, content, \texttt{is\_correct}).
    \item \texttt{exams}: Thông tin kỳ thi (id, title, \texttt{duration\_minutes}, \texttt{start\_time}, \texttt{end\_time}, \texttt{max\_attempts}, password, \texttt{is\_dynamic}, \texttt{dynamic\_config}, status). Cột \texttt{dynamic\_config} lưu cấu hình sinh đề dạng JSONB bao gồm tỷ lệ phân bổ độ khó và số lượng câu hỏi cần lấy theo từng chủ đề.
    \item \texttt{exam\_questions}: Liên kết danh sách câu hỏi cố định trong đề thi.
    \item \texttt{exam\_submissions}: Lưu thông tin bài nộp của học viên (id, \texttt{exam\_id}, \texttt{user\_id}, score, \texttt{started\_at}, \texttt{submitted\_at}).
    \item \texttt{user\_answers}: Lưu câu trả lời chi tiết của thí sinh đối với từng câu hỏi trong lượt làm bài (\texttt{submission\_id}, \texttt{question\_id}, \texttt{chosen\_choice\_id}, \texttt{is\_correct}, \texttt{awarded\_points}).
    \item \texttt{exam\_violations}: Lưu trữ nhật ký vi phạm (id, \texttt{exam\_id}, \texttt{user\_id}, \texttt{violation\_type}, \texttt{violation\_time}) nhằm lưu bằng chứng khi thí sinh rời màn hình hoặc thực hiện sao chép trong phòng thi.
\end{itemize}

\section{Xây dựng ứng dụng}
\label{section:4.3}

\subsection{Thư viện và công cụ sử dụng}
Các công cụ phát triển phần mềm được sử dụng trong đồ án JQK Study được tổng hợp chi tiết trong Bảng \ref{table:my_label}.

\begin{table}[H]
\centering
    \begin{tabular}{lll}
        \hline
        \textbf{Mục đích} & \textbf{Công cụ/Thư viện} & \textbf{Phiên bản/Địa chỉ URL} \\ \hline
        Ngôn ngữ lập trình Backend & Go (Golang) & v1.25.3 / https://go.dev \\ 
        Framework giao tiếp API & gRPC & v1.62.0 / https://grpc.io \\ 
        Bộ thư viện ORM Go & GORM & v1.31.0 / https://gorm.io \\ 
        Cơ sở dữ liệu chính & PostgreSQL & v16.0 / https://www.postgresql.org \\ 
        Hệ thống lưu đệm & Redis & v7.0 / https://redis.io \\ 
        Hệ thống hàng đợi sự kiện & Apache Kafka & v3.5.0 / https://kafka.apache.org \\ 
        Framework phát triển Frontend & Next.js & v16.0.1 / https://nextjs.org \\ 
        Đóng gói ứng dụng & Docker & v24.0 / https://www.docker.com \\ 
        Hệ thống điều phối ứng dụng & Kubernetes & v1.28.2 / https://kubernetes.io \\ 
        Công cụ phát triển hot-reload & Tilt & v0.33 / https://tilt.dev \\ \hline
    \end{tabular}
    \caption{Danh sách thư viện và công cụ phát triển hệ thống}
    \label{table:my_label}
\end{table}

\subsection{Kết quả đạt được}
Hệ thống JQK Study đã được xây dựng và đóng gói hoàn chỉnh thành 5 Docker Container độc lập: `api-gateway`, `user-service`, `exam-service`, `course-service`, và `notification-service`. Tổng số dòng code Go phía backend khoảng 12,000 dòng và Next.js frontend khoảng 8,500 dòng TypeScript. Hệ thống hoạt động trơn tru trên cụm Kubernetes cục bộ được điều khiển bởi Tilt, tự động quản lý kết nối cơ sở dữ liệu và phục vụ giao diện tại cổng 4001.

\subsection{Minh họa các chức năng chính}
Các chức năng chính bao gồm:
1. **Làm bài thi trực tuyến:** Màn hình làm bài được khóa tương tác bên ngoài, đồng hồ tự động đếm ngược và cập nhật đáp án về server mỗi 30 giây (auto-save).
2. **Dashboard giám sát vi phạm:** Giảng viên xem được danh sách sinh viên vi phạm kèm số lần chuyển tab cụ thể và thời gian vi phạm chính xác.
3. **Sinh đề thi tự động:** Giảng viên chọn chủ đề, nhập tỷ lệ độ khó và hệ thống tự động bốc đề ngẫu nhiên từ ngân hàng câu hỏi.

\section{Kiểm thử và Triển khai}
\label{section:4.4}

Hệ thống được kiểm thử thông qua các phương pháp kiểm thử đơn vị (Unit Test) cho các hàm xử lý tính toán điểm thi, xáo trộn câu hỏi và kiểm thử tích hợp (Integration Test) cho luồng nghiệp vụ từ khi bắt đầu thi, tự động lưu câu trả lời, gửi cảnh báo vi phạm cho đến khi nộp bài thi thành công.

Về triển khai, hệ thống JQK Study được cấu hình triển khai trên nền tảng Kubernetes. Trong môi trường kiểm thử phát triển, Tilt được cấu hình để tự động hóa quá trình biên dịch mã nguồn, build lại docker image và deploy nóng lên cụm Kubernetes cục bộ mỗi khi lập trình viên lưu mã nguồn. Khi đưa vào vận hành thực tế, hệ thống có thể được dễ dàng phân phối lên các dịch vụ đám mây (AWS EKS hoặc Google GKE) với khả năng tự động co giãn số lượng Pod dựa trên mức độ sử dụng CPU thông qua cấu hình Horizontal Pod Autoscaler.

\end{document}
